\subsection{射影空间结构}

给定一个集合 E 和一个域 K，E 上关于域 K 的（左）射影空间结构由给出一个非空集合 \(\Phi\) ——它由射影空间 \(\mathbf{P}(\mathbf{K}_{s}^{\mathrm{(E)}})\) 的一些子集到 E 上的双射组成——来定义，且满足以下公理：

\((\mathbf{EP}_{\mathrm{I}})\) 每个映射 \(f \in \Phi\) 的定义集是 \(\mathbf{P}(\mathbf{K}_s^{(\mathbf{E})})\) 的一个线性簇。

\((\mathrm{EP}_{\Pi})\) 对每一对元素 \(f, g\) （属于 \(\Phi\) ，分别定义在线性簇 \(\mathbf{P}(\mathbf{V})\) 与 \(\mathbf{P}(\mathbf{W})\) 上），双射 \(h = \overline{g}^{-1} \circ f\) （ \(\mathbf{P}(\mathbf{V})\) 到 \(\mathbf{P}(\mathbf{W})\) 上）是一个射影映射。

\((\mathbf{EP}_{\mathrm{III}})\) 反之，若 \(f \in \Phi\) 定义在线性簇 \(\mathbf{P}(\mathbf{V})\) 上，而 \(h\) 是 \(\mathbf{P}(\mathbf{V})\) 到一个线性簇 \(\mathbf{P}(\mathbf{W}) \subset \mathbf{P}(\mathbf{K}_s^{\mathbf{(E)}})\) 上的双射射影映射，则 \(f \circ h^{-1} \in \Phi\) 。

设 E 是一个集合， \((\mathbf{V}_{\lambda})_{\lambda \in L}\) 是 K 上的一族向量空间，并且假设对每个 \(\lambda \in L\) 给定了一个双射 \(f_{\lambda}\) （从 \(\mathbf{P}(\mathbf{V}_{\lambda})\) 到 E 上），使得对每一对有序指标 \(\lambda, \mu, f_{\lambda} \circ f_{\mu}\) 而言，它是 \(\mathbf{P}(\mathbf{V}_{\mu})\) 到 \(\mathbf{P}(\mathbf{V}_{\lambda})\) 的一个射影映射。于是我们可以如下在 E 上定义关于 K 的射影空间结构：设 \((e_{t})_{t \in I}\) 是空间 \(V_{\lambda}\) 的一个基，记 \(a_{t} = f_{\lambda}(\pi(e_{t}))\) ；设 \(b_{t}\) 是指标为 \(a_{t}\) 的元素（在 \(\mathbf{K}_{s}^{(\mathbf{E})}\) 的典范基中）（§1，第 11 号）。关系 \(t \neq k\) 蕴含 \(b_{t} \neq b_{k}\) ，这是因为 \(f_{\lambda}\) 是双射这一假设；因此诸 \(b_{t}\) 构成向量子空间 \(W_{0}\) （ \(K_{s}^{(\mathbf{E})}\) 的）的一个基，从而存在 \(\mathbf{P}(\mathbf{W}_{0})\) 到 \(\mathbf{P}(\mathbf{V}_{\lambda})\) 上的一个双射射影映射 h ，使得 \(h(\pi(b_{t})) = \pi(e_{t})\) 对所有 \(t \in I\) 成立。若取 \(\Phi\) 为所有双射射影映射 \(f_{\lambda} \circ h \circ g^{-1}\) 的集合，其中 g 遍历所有双射射影映射 \(\mathbf{P}(\mathbf{W}) \subset \mathbf{P}(\mathbf{K}_{s}^{(\mathbf{E})})\) 的集合，则立即可验证 \(\Phi\) 满足公理 \((\mathrm{EP}_{\mathrm{I}}), (\mathrm{EP}_{\mathrm{II}})\) 与 \((\mathrm{EP}_{\mathrm{III}})\) 。此外显然 \(\Phi\) 既不依赖于指标 \(\lambda \in L\) 的选取，也不依赖于基 \((e_{t})\) （在 \(V_{\lambda}\) 中）的选取，也不依赖于 h 的选取。

特别地（取 L 由单个元素组成），由向量空间 V 导出的每个射影空间 \(\mathbf{P}(\mathbf{V})\) 因此具有本号所给定义意义下的一个完全确定的“射影空间结构”。因此，任何具有射影空间结构的集合都可以称为射影空间。

\[
f \in \Phi
\]

\[
\mathbf {P} (\mathbf {V}) \subset \mathbf {P} \left(\mathrm{K} _ {\mathrm{s}} ^ {(\mathrm{E})}\right)
\]

\(f(\mathbf{M})\) 是 \(\mathbf{P}(\mathbf{V})\) 中第 7 号意义下的线性簇（此性质于是对所有 \(f \in \Phi\) 成立）。由此可知，射影空间中的每个线性簇都典范地具有一个射影空间结构。

称射影空间 \(\mathbf{E}\) 的维数为 \(n\) ，如果对所有 \(f\in\Phi,f(\mathbf{E})\) 都是维数为 \(n\) 的线性簇（只需对一个映射 \(f\in\Phi\) 成立即可）。

